\subsubsection{Vertical boundary in a proper curve fibration}

\begin{proposition}[Proved: the homotopy sequence after deleting vertical boundary]
\label{D:vertical-fibration-exact}
Let $k$ be algebraically closed, let $X$ be a smooth connected
proper surface, and let $f:X\to B$ be a morphism to a smooth
connected proper curve.  Let $B^\circ\subset B$ be a nonempty
open over which $f$ is smooth with geometrically connected
one-dimensional fibres.  Let $U\subset X$ be an open containing
$X^\circ=f^{-1}(B^\circ)$.  Suppose that for every
$b\in B\setminus B^\circ$, the open $U$ contains a point of
$X_b$ at which $f$ is smooth.  Then, with compatible geometric
base points, the sequence
\[
 \pi_1(X_{\overline\eta})\longrightarrow\pi_1(U)
       \longrightarrow\pi_1(B)\longrightarrow1
\]
is exact.  If $g$ is the genus of the geometric generic fibre
and $h$ is the genus of $B$, then $\pi_1(U)$ can be generated
topologically by at most $2g+2h$ elements.
\end{proposition}
\begin{proof}
The schemes $X$, $U$ and $X^\circ$ are normal and integral.
Restriction of a connected finite \'etale cover of $U$ to the
dense open $X^\circ$ remains connected.  Thus
\[
 H:=\pi_1(X^\circ)\longrightarrow G:=\pi_1(U)
\]
is surjective.  The proper smooth homotopy sequence gives
\[
 \pi_1(X_{\overline\eta})\longrightarrow H
       \longrightarrow\pi_1(B^\circ)\longrightarrow1
\]
by \cite[X, Corollary 1.4, p. 263]{SGA1}.  Let $N$ be the image
of the first group in $G$.  It is a closed normal subgroup,
and $\pi_1(B^\circ)$ surjects onto $G/N$.

We show that this surjection factors through $\pi_1(B)$.
Let $Q$ be any finite quotient of $G/N$.  It gives a connected
finite \'etale $Q$-torsor $V\to U$.  Its restriction to
$X^\circ$ comes from a finite \'etale $Q$-torsor
$B'^\circ\to B^\circ$, by the proper homotopy sequence.
Normalize $B$ in $k(B'^\circ)$ to obtain the finite map
$B'\to B$.

Fix $b\in B\setminus B^\circ$ and an $f$-smooth closed point
$x\in U\cap X_b$.  Such a closed point exists in the indicated
nonempty smooth locus.  Over the strictly henselian trait
$T=\operatorname{Spec}\mathcal O^{\mathrm{sh}}_{B,b}$, smoothness
gives a section $s:T\to U\times_B T$ specializing to $x$.
To see this directly, choose \'etale coordinates for $f$ at $x$,
lift their residue values to $T$, and use the section of the
resulting \'etale neighborhood of the closed point of $T$;
see \cite[tags 054L and 03QO]{C_StacksFibration}.
The pullback $s^*V$ is finite \'etale over $T$.  On its generic
point it is precisely the base change of $B'^\circ$.
Since $T$ is normal, this finite \'etale cover is the
normalization of $T$ in that generic algebra.  Thus $B'\to B$
is \'etale at $b$.  Doing this at every omitted base point
shows that $B'\to B$ is finite \'etale.

Every finite quotient of $G/N$ therefore kills the kernel of
$\pi_1(B^\circ)\to\pi_1(B)$.  Their intersection is trivial,
so the surjection to $G/N$ factors continuously through
$\pi_1(B)$.  The natural map $G\to\pi_1(B)$ induced by $f$
kills $N$.  These two maps between $G/N$ and $\pi_1(B)$ are
inverse, as is checked after the common surjection from
$\pi_1(B^\circ)$.  This proves exactness.

A smooth proper curve of genus $r$ over an algebraically closed
field has a fundamental group generated by $2r$ elements
\cite[X, Theorem 2.6, pp. 272--273]{SGA1}.
Thus $N$ has at most $2g$ generators and $G/N$ has at most
$2h$ generators.  A generating set of $N$, together with lifts
of a generating set of $G/N$, generates $G$ topologically.
This gives the stated bound.
\end{proof}

\begin{proposition}[Proved: vertical boundary meeting every fibre]
\label{D:vertical-fibration-multiplicity}
Let $k$ be algebraically closed, let $X$ be a normal connected
proper surface, and let $f:X\to B$ be a morphism to a smooth
connected proper curve.  Suppose that over a nonempty open
$B^\circ\subset B$ the morphism is flat with geometrically
reduced, geometrically connected one-dimensional fibres.
Let $U\subset X$ be an open containing $X^\circ=f^{-1}(B^\circ)$
and meeting every fibre of $f$.  Then $\pi_1(U)$ is topologically
finitely generated.  No restriction on the multiplicities of
the exceptional fibres is required.

More precisely, let $N$ be the image of
$\pi_1(X_{\overline\eta})$ in $\pi_1(U)$ and put
$Q=\pi_1(U)/N$.  For each $b\in B\setminus B^\circ$, let $m_b$
be the gcd of the multiplicities of the fibre components whose
generic points lie in $U$.  The image in $Q$ of the full inertia
group at $b$ is finite of order dividing $m_b$.
If all $m_b$ equal one, then $Q=\pi_1(B)$ and the homotopy
sequence in Proposition~\ref{D:vertical-fibration-exact} is exact.
If the geometric generic fibre is smooth of genus $g$, $B$ has
genus $h$, and every $m_b$ is prime to the positive
characteristic, then $\pi_1(U)$ has at most $2g+2h+r$ generators,
where $r=\#(B\setminus B^\circ)$.
\end{proposition}
\begin{proof}
Since $X$ is normal and integral, restriction to $X^\circ$
gives a surjection $H=\pi_1(X^\circ)\to G=\pi_1(U)$.
SGA~1 uses ``separable'' for a flat morphism with geometrically
reduced fibres.  Its proper homotopy sequence therefore applies
to $f|_{X^\circ}$ exactly under the present assumptions
\cite[X, Definition 1.1 and Corollary 1.4, pp. 261--263]{SGA1}.
Thus $N$ is closed and normal, and there is a surjection
\[
                    \pi_1(B^\circ)\longrightarrow Q.
\]
The proper reduced connected geometric generic curve has a
finitely generated fundamental group: normalize its finitely
many irreducible components, use the proper smooth-curve theorem
on each normalization, and apply finite proper descent
\cite[IX, Corollary 5.2; X, Theorem 2.6]{SGA1}.
Consequently $N$ is finitely generated.

Let $Q_0$ be any finite quotient of $Q$.  As in the preceding
proof, the corresponding finite \'etale cover $V\to U$ is over
$X^\circ$ the pullback of a connected finite Galois cover
$B'^\circ\to B^\circ$ with group $Q_0$.  Normalize $B$ in its
function field to obtain $B'\to B$.

Choose a fibre component $D$ over $b$ whose generic point lies
in $U$, and let $m$ be its multiplicity.  Such a component
exists because $U$ meets the fibre and its intersection with
the fibre is open.  Choose a prolongation to $k(V)$ of the
$D$-valuation of $k(X)$.  The first ramification index is one,
since the cover is \'etale on $U$.  The restriction of the
$D$-valuation to $k(B)$ has index $m$.  On the other hand,
$k(B')$ is an intermediate field of $k(V)/k(B)$.  Multiplication
of indices of these discrete value groups yields
\[
                         e(b'/b)\mid m.
\]
This argument applies to every retained component.  Thus
$e(b'/b)$ divides their gcd $m_b$.

Every closed point of $B'$ has residue field $k$.  Since the
base cover is Galois, its full inertia group at $b'$ has order
$e(b'/b)$; here finite extensions of the excellent curve DVRs
satisfy the usual degree formula.  It follows that in every
finite quotient of $Q$, the image of inertia at $b$ has order
dividing $m_b$.  Its image $I_b$ in $Q$ is therefore finite,
with order dividing $m_b$.  Indeed, the orders of all its finite
quotients are bounded; choose a quotient of maximal order and
observe that every refinement has the same kernel on $I_b$.
The intersection of these kernels is trivial.

There are only finitely many omitted base points.  Choose an
open normal subgroup $Q'\subset Q$ meeting every $I_b$ trivially;
this is possible by separating the finitely many nonidentity
elements of these finite groups in finite quotients of $Q$.
Let $B_1^\circ\to B^\circ$ be the connected finite \'etale cover
corresponding to the inverse image of $Q'$ in $\pi_1(B^\circ)$,
and let $B_1$ be its smooth proper compactification.  Every
inertia group of $B_1^\circ$ above a point of
$B\setminus B^\circ$ is an intersection with a conjugate of the
corresponding base inertia.  Its image in $Q'$ is therefore
trivial.  Hence the surjection
\[
                  \pi_1(B_1^\circ)\longrightarrow Q'
\]
factors through $\pi_1(B_1)$: a finite cover of a smooth open
curve extends \'etale to the proper curve exactly when all
these inertia groups act trivially.  The proper smooth-curve
theorem makes $Q'$ finitely generated.  A profinite group with
a finitely generated open subgroup is finitely generated, so
$Q$ is finitely generated.  Since $N$ is finitely generated,
so is $G$.

If $m_b=1$ at every omitted point, all $I_b$ are already
trivial.  The argument in
Proposition~\ref{D:vertical-fibration-exact} then identifies
$Q$ with $\pi_1(B)$ and proves the asserted exactness.
If instead all $m_b$ are prime to the characteristic, every
finite base cover just considered is tame at infinity.
Thus $Q$ is a quotient of $\pi_1^{\mathrm t}(B^\circ)$, which
has at most $2h+r$ generators
\cite[XIII, Corollary 2.12, pp. 392--393]{SGA1}.
When the geometric generic fibre is smooth, $N$ has at most
$2g$ generators, giving the stated numerical bound.
\end{proof}

\begin{theorem}[Proved: finite punctures on a proper surface mapping to a curve]
\label{D:surface-map-to-curve}
Let $S$ be a normal connected proper surface over an
algebraically closed field $k$.  If $S$ admits a nonconstant
morphism to a proper curve, then
\[
                         \pi_1(S\setminus Y)
\]
is topologically finitely generated for every finite closed
subset $Y\subset S$.
The morphism need not be generically smooth: geometrically
reduced singular generic fibres are allowed.
\end{theorem}
\begin{proof}
Replace the target by its reduced image and take Stein
factorization.  We obtain a proper surjective map $f:S\to B$
with $f_*\mathcal O_S=\mathcal O_B$ and geometrically connected
fibres.  The curve $B$ is normal, hence smooth and proper over
$k$, and $F=k(B)$ is algebraically closed in $K=k(S)$.
These assertions follow from Stein factorization and normality
of $S$ \cite[Theorem 37.53.4, tag 03H0]{C_StacksFibration}.
To make the constant-field assertion explicit, the global
functions on the proper normal generic fibre are exactly the
relative algebraic closure of the base field in $K$: an element
algebraic over the base is integral over every local ring of
that fibre and hence is regular there by normality.

We check generic separability in positive characteristic.
The field $F$ has transcendence degree one over the perfect
field $k$, so $[F:F^p]=p$.  Choose $t\in F\setminus F^p$;
then $F^{1/p}=F(t^{1/p})$.  A $p$th root of $t$ in $K$ would
be algebraic over $F$ and therefore lie in $F$, contradicting
the choice of $t$.  It follows that
\[
 K\otimes_F F^{1/p}\simeq K[T]/(T^p-t)
\]
is a field.  The separability criterion
\cite[Lemma 10.44.2, tag 05DT]{C_StacksFibration} implies that
$K/F$ is separable and geometrically reduced.  In characteristic
zero this conclusion is automatic.  Relative algebraic
closedness also gives geometric irreducibility
\cite[Lemma 10.47.8, tag 037P]{C_StacksFibration}.
Thus the generic fibre is geometrically integral.  This does
not assert that it is smooth.

The morphism $f$ is flat: at a point above a closed point of
$B$, its local ring is torsion-free over the DVR of $B$, hence
flat.  Geometric reducedness of the generic fibre spreads to
a nonempty open of $B$, by the openness theorem for flat proper
morphisms
\cite[Lemma 37.26.7, tag 0C0E]{C_StacksFibration}.
Shrink this open further to avoid $f(Y)$ and denote it by
$B^\circ$.  Its fibres are geometrically reduced and connected,
so the proper separable homotopy theorem applies.

Put $U=S\setminus Y$.  Then $f^{-1}(B^\circ)\subset U$.
Every fibre of the dominant map from the integral surface to
the smooth curve has a one-dimensional component; equivalently,
the local fibre parameter cuts a nonzero principal divisor.
Deleting finitely many closed points leaves the generic points
of all these components.  Thus $U$ meets every fibre of $f$.
Proposition~\ref{D:vertical-fibration-multiplicity} gives the
claim.
\end{proof}

\begin{corollary}[Proved: birational and finite-normalization descent of the fibration criterion]
\label{D:vertical-fibration-models}
Let $S$ be a normal connected proper surface over an
algebraically closed field, and let $V\subset S$ be a nonempty
open.  Suppose a proper birational map $q:X\to S$ from a
smooth proper surface has a curve fibration for which
$U=q^{-1}(V)$ satisfies either
Proposition~\ref{D:vertical-fibration-exact} or
Proposition~\ref{D:vertical-fibration-multiplicity}.
Then $\pi_1(V)$ is topologically finitely generated, and inherits any generator bound supplied by the relevant case.  In particular, the same is
true of any fixed-boundary tame quotient of $\pi_1(V)$.

For a possibly nonnormal proper scheme of dimension at most
two, the same conclusion holds if its reduced components have
finite normalizations and their inverse-image opens satisfy
this condition, with lower-dimensional components having
finitely generated fundamental groups.  In particular, the
latter requirement is automatic for a boundary of codimension
at least two in every component.
\end{corollary}
\begin{proof}
A connected finite \'etale cover of the normal integral $V$
has a single function field.  Pulling it back along the
birational map $U\to V$ preserves that generic field.  Every
finite \'etale component over the normal integral $U$
dominates $U$, so the pullback remains connected.  Hence
$\pi_1(U)\to\pi_1(V)$ is surjective.  The first assertion and any stated generator bound follow from the propositions.

For the second assertion use the finite proper surjection
from the disjoint union of the normalizations of the reduced
components, restricted to the indicated opens.  Finite proper
descent of finite \'etale covers, in the form already used in
Proposition~\ref{D:fixed-boundary-positive}, shows that finite
generation descends
\cite[IX, Corollary 5.2]{SGA1}.  Under the final codimension
condition, components of dimension at most one lose no points
and remain proper, so their fundamental groups are finitely
generated.
\end{proof}
